% DONO: João Pedro (exemplo, construções) + Nathan (decisões)   ·   ESPAÇO: ~1,25 página
\section{A linguagem}
% O enunciado pede: programa de exemplo, as construções principais e TRÊS
% decisões de projeto justificadas, cada uma com a alternativa descartada.

\subsection{Exemplo}
% Usar exemplos/validos/promocao_outubro.vitrine e o recibo. Explicar em
% prosa: produto, regra quando/entao, carrinho, calcular.
\begin{lstlisting}
produto camiseta = R$ 49,90;
produto caneca   = R$ 29,90;
frete = R$ 15,00;
regra leve3 { quando qtd(camiseta) >= 3
              entao desconto 20% em camiseta; }
regra frete_gratis { quando total >= R$ 150,00
                     entao frete R$ 0,00; }
carrinho { camiseta x 3; caneca x 1; }
calcular;
\end{lstlisting}

\subsection{Construções e tipos}
% - Tipos: inteiro, dinheiro (centavos inteiros), percentual (só em desconto).
% - Precedência pela cadeia expressao > termo > fator. Ex.:
%   R$ 10,00 + R$ 2,00 * 3 = R$ 16,00.
% - Escopo global único; declarar antes de usar.
% - Uma tabela pequena de tipos (docs/especificacao.md, seção 5).

\subsection{Decisões de projeto}
% Três decisões (guia, seção 6). Para cada uma: decisão, alternativa, por quê.
% 1. Descontos aditivos sobre o preço original x descontos compostos.
%    SEJAM HONESTOS: com composição a soma nunca passa de 100%, então a
%    escolha cria o invariante que a verificação protege. Digam por que
%    ainda assim faz sentido (previsibilidade para o lojista).
% 2. total = subtotal antes dos descontos x total corrente na ordem das
%    regras (que tornaria o resultado dependente da ordem e invalidaria a
%    análise do conjunto). Caso concreto: no exemplo canônico, R$ 179,60
%    antes do desconto e R$ 149,66 depois, com frete grátis a partir de
%    R$ 150,00.
% 3. Análise conservadora por faixas x soma ingênua (recusa faixas
%    exclusivas, o caso mais comum) x solver exato (SMT: biblioteca pronta,
%    proibida pelo enunciado para a semântica).
